Все 48 задач основной серии прошли проверку полноты, хешей завершения и исходных траекторий. Ни один из 528 калибровочных и валидационных эпизодов не достигает критерия тяжёлого события. Максимальный WIP при калибровке равен 16, 61, 19 и 111 для фиксированного мира, обслуживания, возвратов и совместного механизма соответственно. Последнее значение включает краткое превышение 100, поэтому <<до кризиса>> означает отсутствие заданного длительного события, а не постоянно низкую нагрузку. Ни один эпизод не исключён.

Таблица~\ref{tab:e7cost} показывает все механизмы и методы сравнения. В основном условии с возвратами планирование с учётом нагрузки выбирает сокращение возвратов при пороге 20 в 9 из 12 обучений, добавление мощности при пороге 20 --- в остальных 3. Средняя фактическая стоимость снижается с 92,832.6 у статического планирования до 57,351.8, то есть на 38.2\%. Сравнение с добавлением мощности при пороге 20 даёт снижение на 38.1\%. Оба скорректированных интервала в табл.~\ref{tab:e7contrast} исключают ноль: зафиксированный основной критерий выполнен. Фиксированное сокращение возвратов при пороге 20 даёт стоимость 45,313.5: результат модели численно хуже, а парный интервал включает ноль. Это не доказывает ни равенства, ни превосходства над данным правилом. Выбор компонент по прогнозной точности сохраняет зависимость возвратов от нагрузки лишь в 2 из 12 обучений и выбирает сокращение возвратов лишь в 1; его средняя стоимость равна 88,834.8.

\begin{table}[t]
\caption{Средняя фактическая стоимость E7 в тысячах условных единиц времени задачи, по 12 независимым обучающим наборам на механизм и всем тестовым условиям. <<Отбор>> означает выбор компонент по валидации.}\label{tab:e7cost}
\centering\small\setlength{\tabcolsep}{3pt}
\begin{tabular}{lrrrrrr}\toprule
Механизм & \shortstack{Без\\действия} & \shortstack{Статич.\\модель} & \shortstack{С\\нагрузкой} & \shortstack{С\\отбором} & \shortstack{Мощность\\при 20} & \shortstack{Возвраты\\при 20}\\\midrule
Фиксированный & 23.6 & 18.0 & 19.9 & 18.0 & 17.7 & 22.6 \\
Обслуживание & 211.7 & 111.0 & 111.0 & 111.0 & 111.0 & 202.7 \\
Возвраты & 147.8 & 92.8 & 57.4 & 88.8 & 92.6 & 45.3 \\
Совместный & 375.2 & 309.2 & 288.3 & 292.5 & 309.2 & 284.1 \\
\bottomrule

\end{tabular}
\end{table}
\begin{table}[t]
\caption{Возвраты: стоимость планирования с нагрузкой минус метод сравнения, в тысячах. Основные интервалы имеют уровень 97.5\%; сравнение с сокращением возвратов описательное, 95\%. Отрицательная разность означает выигрыш модели с нагрузкой.}\label{tab:e7contrast}
\centering\small
\begin{tabular}{lrrr}\toprule
Метод сравнения & Разность & Парный интервал & Уровень\\\midrule
Статич. планирование & $-35.48$ & $[-51.56, -19.40]$ & 97.5\% \\
Мощность при 20 & $-35.27$ & $[-51.23, -19.32]$ & 97.5\% \\
Возвраты при 20 & $12.04$ & $[-1.87, 25.94]$ & 95\% \\
\bottomrule

\end{tabular}
\end{table}

Реакция включается раньше кризиса в парной траектории без действия во всех 289 опасных траекториях мира с возвратами; средний запас времени равен 108.0 единицы. Она предотвращает заданное событие в 93 траекториях; в 196 оно сохраняется. В 192 траекториях без всплеска нет ни срабатываний, ни кризисов. Рис.~\ref{fig:e7} показывает ускоренное восстановление и разброс, вызванный различиями выбранных правил. Это условные описательные измерения предупреждения, а не оценка калиброванного классификатора кризисов.

\begin{figure}[t]
\centering\includegraphics[width=\textwidth]{generated/e7_response_ru.pdf}
\caption{Реакции в мире с возвратами для двух отложенных сценариев со всплеском. Восемь траекторий усредняются внутри каждого из 12 обучений, затем усредняются обучения. Полоса модели с нагрузкой --- одно стандартное отклонение средних по обучениям, а не доверительный интервал. Серые области отмечают всплески; горизонталь --- WIP 100. Фиксированное правило сокращения возвратов показывает ограничение выбора действия по модели.}\label{fig:e7}
\end{figure}

Польза зависит от механизма. При влиянии нагрузки на обслуживание все планировщики выбирают добавление мощности при пороге 20 и дают одинаковую стоимость. При совместном механизме модель с нагрузкой снижает стоимость на 6.7\% относительно статического планирования и правила добавления мощности, однако частота тяжёлых событий растёт с 69.3\% до 71.0\%: минимизация стоимости очереди не равна предотвращению заданного события. Без обратной связи модель с нагрузкой увеличивает стоимость на 10.9\% относительно статической, с описательной парной разностью 1,959.3 и интервалом $[219.5,3699.1]$. По четырём механизмам частота срабатываний без всплеска равна 0.0\%, 14.1\%, 0.0\% и 28.6\%; соответствующая частота кризисов без действий --- 0.0\%, 2.1\%, 0.0\% и 19.3\%. Неустойчивость базовой работы объясняет часть, но не все срабатывания без всплеска.

Анализ с равными весами механизмов при исходной стоимости даёт выигрыш модели с нагрузкой: нормированные разности равны $-0.0536$ относительно статического планирования, $-0.0503$ относительно добавления мощности при пороге 20 и $-0.1137$ относительно сокращения возвратов при пороге 20. Это разности отношений стоимости внутри миров, а не проценты снижения суммарной ненормированной стоимости. Таблица~\ref{tab:e7sensitivity} показывает описательные интервалы и все заранее заданные множители стоимости. При удвоенной стоимости действий интервалы относительно статического планирования и добавления мощности включают ноль; выигрыш относительно сокращения возвратов остаётся пограничным. Таким образом, приспособление реакции к механизму подтверждено для заданной задачи принятия решения, а не независимо от предположений о стоимости.

\begin{table}[t]
\caption{Нормированные разности стоимости по механизмам: планирование с нагрузкой минус метод сравнения, с описательными 95\%-ными интервалами. Для каждого множителя правила заново выбираются по стоимости планирования.}\label{tab:e7sensitivity}
\centering\small\setlength{\tabcolsep}{3pt}
\begin{tabular}{lrrr}\toprule
Множитель & Статич. планирование & Мощность при 20 & Возвраты при 20\\\midrule
0.5 & \shortstack{$-0.0467$\\$[-0.0772,-0.0161]$} & \shortstack{$-0.0467$\\$[-0.0772,-0.0161]$} & \shortstack{$-0.1201$\\$[-0.1501,-0.0901]$} \\
1.0 & \shortstack{$-0.0536$\\$[-0.0824,-0.0247]$} & \shortstack{$-0.0503$\\$[-0.0790,-0.0217]$} & \shortstack{$-0.1137$\\$[-0.1420,-0.0855]$} \\
2.0 & \shortstack{$-0.0136$\\$[-0.0618,0.0346]$} & \shortstack{$-0.0037$\\$[-0.0510,0.0436]$} & \shortstack{$-0.0471$\\$[-0.0939,-0.0002]$} \\
\bottomrule

\end{tabular}
\end{table}
